\section{The Person under Overlapping Orders: An Obligation Account}

The several orders of law meet in the action of one morally responsible person. Their unity is not a super-code or a single human appellate hierarchy. Natural and revealed norms measure the morality of action; civil and ecclesiastical orders establish juridical effects within their competence; offices, vows, contracts, family relations, and other commitments add role-specific duties. Conscience judges the concrete act, but it does not legislate, dispense, or adjudicate for itself.

\subsection{First establish the person, act, and facts}

An obligation account begins with the act to be done or omitted, not with an institutional label. Identify the authoritative wording, material facts, relevant time, the person's actual role, the conduct required, the persons affected, and the foreseeable consequences. A duty to report, a prohibition on disclosure, a sacramental decision, and an employment command may concern the same events while addressing different acts and effects.

\begin{longtable}{>{\bfseries\raggedright\arraybackslash}p{0.20\linewidth}>{\raggedright\arraybackslash}p{0.34\linewidth}>{\raggedright\arraybackslash}p{0.38\linewidth}}
\toprule
\textbf{Relation to audit} & \textbf{Facts to establish} & \textbf{Why the facts matter}\\
\midrule
\endhead
Moral agent & Knowledge, freedom, intention, object of the act, circumstances, alternatives, effects on others, and prior commitments. & Natural moral law has universal objective reach, while responsibility for an act also depends on knowledge and freedom.\\
Canonical person & Baptism or reception; Latin or Eastern ascription; age and use of reason; domicile, quasi-domicile, presence, traveler or transient status; and proper pastor or ordinary. & CIC cc.~96--112 and CCEO cc.~7, 27--38, 909--919 connect canonical personality and authority to these facts. Eastern ascription does not change merely because an Eastern Catholic is entrusted to a pastor of another Church \emph{sui iuris} (CCEO c.~38).\\
Ecclesial condition & Lay, clerical, consecrated, officeholding, institute or association membership, delegation, mandate, vow, sacramental status, and applicable universal, particular, or proper law. & A duty can arise from a condition or validly assumed role even when it does not bind every Catholic. The instrument creating the role and any limits on competence must be verified.\\
Civil relation & Citizenship, residence, physical presence, age and capacity, family status, license, employment, office, fiduciary role, property, corporate form, and pending process. & Civil jurisdiction and effects are state-specific. Canonical status does not answer civil capacity, liability, reporting, employment, or procedural questions.\\
Transaction or act & Contract, decree, precept, rescript, privilege, dispensation, judgment, notice, consent, form, and factual trigger. & General law, a singular administrative act, a contract, and a judgment become operative in different ways and can address different persons.\\
\bottomrule
\end{longtable}

The status map is descriptive, not determinative. It shows which sources must be checked. A person may occupy several roles at once, and an authority competent over one role may lack competence over another.

\subsection{Construct a norm ledger}

List each potentially applicable norm on its own row. Do not combine a moral principle, a statute, a canon, and an employment instruction into one paraphrase.

\begin{fourstable}{Norm and exact source}{Claimed subjects and scope}{Claimed effect}{Authority, forum, and relief}
Natural moral norm & Every person or a class defined by the nature of the act or relationship. & Objective moral duty or right; culpability still requires a personal analysis. & Reasoned moral judgment, authoritative moral teaching where applicable, and competent pastoral or theological counsel; no human dispensation from the underlying norm.\\
Revealed divine norm & Addressees depend on the proposition: all persons, a covenantal people, the baptized, an officeholder, or persons in a sacramental relation. & Moral, doctrinal, sacramental, or constitutional consequence as the source establishes. & The claimed revelation, authoritative teaching, and competent ecclesiastical application must be identified.\\
Ecclesiastical positive norm & Persons defined by baptism or reception, Church \emph{sui iuris}, territory, personal status, office, institute, or a singular act. & External or internal canonical duty, validity, liceity, status, penalty, presumption, or remedy. & Legislator, executive authority, tribunal, or law-governed internal forum; exception, privilege, dispensation, recourse, or appeal only as law provides.\\
Civil positive norm & Persons, conduct, territory, property, entity, or legal relation within the civil authority's jurisdiction. & Civil validity, prohibition, permission, liability, penalty, status, evidence, or remedy. & Competent agency or court under the applicable constitution, statutes, regulations, conflicts rules, and procedure.\\
Assumed or role-based norm & A person bound by contract, vow, office, delegation, professional rule, fiduciary relation, or family responsibility. & The specific performance, loyalty, care, confidentiality, reporting, or other duty established by the relation. & The instrument, governing law, competent superior or counterparty, and any procedure for release, modification, or adjudication.\\
\end{fourstable}

\subsection{Apply the binding-force and manifestation gates}

For each positive norm, verify:

\begin{enumerate}
\item the authority's competence over the person, subject matter, territory, and juridical form;
\item the norm's rank and relation to higher, special, particular, proper, and later law within the same order;
\item the authoritative text, promulgation and effective date for a general law, or legitimate notification and execution for a singular act;
\item the personal, territorial, temporal, and factual conditions that bring the norm into operation;
\item the current interpretation, including any authentic interpretation or controlling judgment;
\item any exception, privilege, acquired right, dispensation, excuse, necessity rule, or doubt of law or fact;
\item the effect of ignorance or error under the rule governing that particular consequence; and
\item whether a treaty, concordat, reception rule, contract, or recognition rule connects the norm to another legal order.
\end{enumerate}

Latin canon law illustrates the distinctions. CIC cc.~7--22 govern promulgation, subjects, territorial scope, doubt, ignorance, interpretation, later law, and reception of civil law. A singular precept under c.~49 directly enjoins a determined person, while a singular decree ordinarily has effect through legitimate notification under cc.~52--56. Dispensation concerns a merely ecclesiastical law and requires competent authority and a just cause (cc.~85--93). The CCEO supplies its own corresponding rules in cc.~1488--1504 and 1510--1539. An apparent permission in one order removes only the obstacle addressed by that order; it does not establish the act's morality or erase a duty in another order.

\subsection{Keep an effects ledger}

The same facts can produce different answers in different columns:

\begin{threestable}{Question}{Required conclusion}{Frequent error}
Moral object and obligation & What act is this person proposing to choose, and is it required, permitted, or forbidden under the applicable natural and revealed norms? & Treating civil or canonical permission as moral approval.\\
Validity and status & Did a juridical or sacramental act produce the effect claimed in the relevant order? & Treating reduced culpability or a good intention as curing invalidity.\\
External obligation and enforceability & Does a civil, canonical, contractual, or administrative norm bind externally, and in which forum can it be enforced? & Assuming moral non-obligation automatically removes every external effect.\\
Liceity and discipline & If the act was valid, was it also licit, and what disciplinary consequence follows from any illicit conduct? & Treating validity, liceity, permission, and impunity as synonyms.\\
Culpability and sanction & What knowledge, freedom, intent, negligence, excuse, presumption, or defense governs moral fault and each penal order? & Importing a moral conclusion into penal liability, or importing a penal excuse into objective morality.\\
Evidence and remedy & What must be proved, by whom, before which authority, by what deadline, and with what relief or suspensive effect? & Treating a correct substantive claim as self-executing.\\
\end{threestable}

Under CIC c.~1321, the accused is presumed innocent until the contrary is proved, and punishment requires an externally committed violation gravely imputable by malice or culpability. Once the external violation has been committed, imputability is presumed unless otherwise apparent. Canons 1322--1325 govern excuses and mitigation. CCEO c.~1414 uses a distinct scheme: an external violation of a penal law or precept is presumed deliberate unless the contrary is proved, while violation of another law or precept is punished only upon repetition after a penal warning. Canons 1415--1416 govern further excuses and mitigation. Moral cooperation, canonical participation in a delict, and civil accomplice liability are different analyses even when they concern the same conduct.

\subsection{Conscience and the several forums}

Conscience is the reasoned judgment that this concrete act should be done or avoided (CCC 1777--1793). A person must form conscience and follow a judgment held without practical doubt. Subjective certainty does not establish the truth of every premise, supply a legal dispensation, or decide external rights. Culpable and nonculpable error can affect personal responsibility differently while leaving objective harm, invalidity, civil consequences, or duties of repair in place.

The canonical internal forum is also not another name for conscience or private advice. CIC c.~130 provides that power of governance is ordinarily exercised for the external forum and sometimes for the internal forum, with external effects recognized only as law provides. CCEO c.~980 distinguishes external from sacramental and nonsacramental internal fora in similar terms. A confessor or adviser ordinarily cannot alter external canonical status without a specific faculty and a legal rule giving the act that effect. The civil forum remains separate.

\subsection{Resolve apparent conflicts before accepting a true conflict}

Concurrent duties often can be harmonized, sequenced, delegated, or accommodated. Within one legal order, interpretation should first consider text, context, purpose, higher rank, special law, later law, and the governing conflict rules. Across civil and ecclesiastical orders, neither forum is universally superior to the other as a human institution; each determines effects within its competence, subject to the common moral measure.

If the duties remain incompatible:

\begin{enumerate}
\item A person may never choose an intrinsically evil act or formally cooperate in evil. When a well-formed and certain judgment of conscience concludes that a command requires such an act, refusal is morally required. Civil or canonical consequences must still be anticipated and addressed.
\item Other participation in another's wrongdoing requires analysis of the person's own act and intention, causal proximity, alternatives, duress, proportionate reason, scandal, and effects on vulnerable third parties. \emph{Evangelium vitae} 73--74 applies these distinctions specifically to laws attacking human life; its conclusions cannot simply be transferred to every dispute.
\item When two morally permissible affirmative duties cannot both be performed, practical reason must weigh their gravity and immediacy, the rights and vulnerability of others, nondelegable role obligations, prior commitments, available substitutes, and the common good. No taxonomic label decides that prudential judgment by itself.
\item Even under an abusive command, residual duties remain: protect persons, avoid needless disorder, provide notice where required, preserve continuity of essential care, repair attributable harm, and render whatever the objective common good still requires (CCC 2238--2242).
\end{enumerate}

\subsection{Use the competent forum and record the decision}

Identify who can give the needed relief: a legislator or authentic interpreter, a competent executive superior or dispenser, a canonical internal or external forum, a tribunal, a civil agency or court, or the relevant contracting authority. Pastoral and legal counsel can assist without possessing adjudicative power. Verify deadlines, exhaustion requirements, available advocates, and whether filing suspends the act. Under the general Latin administrative-recourse sequence, for example, CIC cc.~1734 and 1737 ordinarily provide short periods measured in useful days, and c.~1736 shows that suspension is not generally automatic; CCEO cc.~999--1001 contain an Eastern sequence. Special law can supply different periods and effects.

For a consequential decision, retain a concise account of the controlling text and facts, statuses and roles, norms in each order, interpretation and relief sought, moral judgment, action chosen, persons affected, safeguards, unresolved uncertainty, and the event that should trigger review. This record is not a private decree. It makes the person's reasoning accountable to truth, authority, and the rights of others.
